\section{이론적 적용 범위와 주장하지 않는 사항}

본 방법은 다음의 정확한 성질을 보존한다.
\begin{enumerate}
  \item \textbf{이산 상태 호환성.} 모델은 항상 정수 토큰 인덱스를 입력받는다.
  \item \textbf{첫 스텝의 동치성.} $\alpha_0=0$이므로 같은 난수 추출을 사용할 때 첫 범주형 전이는
  대응하는 기준 전이와 동일하다.
  \item \textbf{단체 유효성.} 사영할 양의 부분이 0이 아니라는 조건 아래 모든 외삽 분포는 음이 아닌
  값으로 정규화된다.
  \item \textbf{MDLM 흡수 구조.} 이미 마스크가 해제된 토큰은 캐시를 사용하지 않는 MDLM 역방향
  전이와 정확히 동일하게 고정된다.
\end{enumerate}

그러나 샘플 시퀀스 $x_{t_i}$는 확률적으로 변하므로 $z_i$는 결정론적 주변분포 곡선 $p_t$가 아니라
무작위 경로를 따라 얻은 조건부 예측들의 모음이다. SEDD에서 $z_i$는 확률흐름 ODE의 상태가 아니라
한 스텝 전이 분포다. 또한 $p_k=0$인 경계에서 Fisher--Rao 계량은 특이하지만, MDLM은 의도적으로
마스크 확률을 0으로 설정하고 이미 마스크가 해제된 토큰에는 델타 분포를 만든다. 제곱근 표현은 이
경계에서도 수치적으로 사용할 수 있으나, 매끄러운 다양체에 기반한 Taylor 가정은 성립하지 않는다.

RDLM은 초구 위에 연속적인 브리지 혼합 SDE를 구성하고 연속 상태를 조건으로 하는 모델을 학습한다
\citep{jo2025rdlm}. Fisher-Flow는 구면 측지선을 따라 확률질량을 운반하는 연속 벡터장을 학습한다
\citep{davis2024fisherflow}. 본 방법은 이 두 작업을 수행하지 않는다. 마찬가지로 DPM-Solver는
확산 ODE의 정확한 준선형 적분 형태를 활용하여 고차 정확도를 얻지만 \citep{lu2022dpmsolver}, 현재
이산 샘플러에는 그와 같은 구조가 없다. 따라서 본 연구는 주변분포 보존, 국소 절단오차 $O(h^3)$,
또는 전역 수렴 $O(h^2)$를 주장하지 않는다.
